\chapter{Introduction}

\section{Problem Importance}
Modern digital environments rely on multiple independent services and applications to perform everyday operations such as monitoring, data synchronization, notifications, and reporting. These operations are often repetitive and require frequent execution, either in response to specific events or according to fixed schedules. Performing such tasks manually consumes significant time, increases operational overhead, and introduces the possibility of human error.

Workflow automation has therefore become an essential approach to improving efficiency and reliability in modern systems. Automation tools enable users to define sequences of tasks that execute automatically, reducing manual intervention and allowing systems to operate continuously. The importance of automation grows as systems become more interconnected and data-driven, making flexible and scalable workflow orchestration a critical requirement in many domains.

\section{Challenges}
Despite the growing adoption of automation platforms, several challenges remain. First, integrating multiple services with different interfaces and data formats introduces complexity in workflow design and execution. Second, supporting both event-driven and periodic execution requires reliable scheduling mechanisms that can operate autonomously over long periods. Third, many existing solutions are either complex to configure or difficult to extend, limiting their suitability for educational or customizable environments.

Another challenge lies in designing a modular architecture that supports future expansion without affecting core system stability. As automation systems evolve, there is increasing demand to incorporate intelligent processing and interaction with external devices, which requires careful architectural planning from the beginning.

\section{Domain}
This project belongs to the domain of workflow automation and orchestration systems. Workflow automation focuses on defining, executing, and monitoring sequences of tasks that connect different services and processes. The project adopts a node-based workflow paradigm in which each node represents a specific operation or integration step.

The domain intersects with distributed systems, event-driven architecture, and task scheduling. Additionally, the system is designed with extensibility in mind, allowing future integration with emerging domains such as artificial intelligence services and Internet of Things (IoT) devices.

\section{Proposed Solution}
The proposed solution is a workflow automation platform inspired by node-based automation tools. The system allows users to create workflows composed of interconnected nodes that execute tasks automatically based on triggers or scheduled intervals.

The platform focuses on providing a clear separation between workflow definition, execution logic, and integration components, ensuring modularity and maintainability. Core features include event handling, periodic task scheduling, workflow execution management, and monitoring capabilities.

While automation and scheduling form the core functionality of the system, the architecture is intentionally designed to support optional extensions such as AI nodes for intelligent processing and IoT nodes for device interaction. These extensions demonstrate the flexibility of the proposed architecture without redefining the core objectives of the project.

\section{Limitations}
The scope of this project is limited to implementing the fundamental concepts of workflow automation and orchestration. The system is intended as an academic and practical prototype rather than a full commercial platform.

Therefore, certain advanced capabilities are out of the scope of this work, including large-scale distributed execution, enterprise-grade security policies, advanced access control mechanisms, and extensive third-party integration marketplaces. AI and IoT components are included only as extensibility demonstrations and are not required for core system operation.

\section{Target Users}
The primary target users of the proposed system are students, developers, and small teams who require a flexible and understandable automation platform for integrating services and executing repetitive tasks. The system is designed to support users who may not wish to build complex automation pipelines manually but still need customization and control over workflow behavior.

By focusing on usability and extensibility, the project aims to provide a learning-friendly automation environment that can serve both educational and practical experimentation purposes.
